\chapter{The Problem and Its Background}

This chapter highlights the critical challenges currently hindering effective
disaster management in Bulacan and CAMANAVA, specifically focusing on the
inefficiency of traditional response mechanisms in processing the massive
influx of unstructured digital communications during disasters and
emergencies. This provides the contextual background and foundation for
implementing this computational solution within the local disaster response
infrastructure.

\section*{Background of the Study}

During severe hydrometeorological events, local disaster management offices
are consistently overwhelmed by a massive influx of unorganized, high-velocity
digital communications from affected citizens. While these localized text
feeds contain critical, actionable information regarding immediate rescue
needs and structural damage, the sheer volume and unstructured nature of the
data render manual processing highly inefficient. To address this information
bottleneck, modern computational semantics relies on vector space mechanics to
mathematically transform qualitative narrative posts into numeric matrices. As
established by \parencite{mihunov2022}, mapping unstructured disaster text into a
high-dimensional geometric space allows automated systems to systematically
extract abstract semantic relationships. This mathematical transformation
reduces chaotic crisis communications into distinct spatial coordinates,
enabling the systematic and automated categorization of emergency surveillance
data.

Building upon this developments, recent computational research
has successfully deployed probabilistic generative models to filter, group,
and route these high-volume disaster reports in real time. Previous empirical
studies demonstrate that algorithms such as Latent Dirichlet Allocation (LDA)
effectively manage the linguistic noise, broken syntax, and data sparsity
inherent in online text feeds during acute natural hazards. For instance,
during the 2023 Hebei Floods, \parencite{zang2026} utilized LDA topic modeling to
process unorganized short-text feeds, proving that probabilistic topic
distributions can successfully categorize heterogeneous public discourse into
actionable, distinct thematic clusters. By treating each text report as a
probabilistic mixture of topics rather than relying on rigid
string-matching, these models allow crisis managers to bypass information
overload and automatically isolate urgent, localized operational needs from
general baseline noise.

While these computational frameworks have demonstrated profound efficacy in
international applications, their deployment is most critical in regions
facing compounding, high-frequency environmental crises. The necessity for
automated, real-time sensemaking tools becomes acutely apparent when
contextualizing the hyper-vulnerable geographic landscape of the Philippine
archipelago. In an environment where municipal emergency management offices
are repeatedly overwhelmed by simultaneous localized crises, the rapid
algorithmic extraction of actionable disaster data is not merely a logistical
optimization, but a fundamental requirement for survival and resilience. This
operational urgency is directly driven by the country's severe climatic
reality.

Disaster-related risks in the Philippines are numerous, as the country is
heavily exposed to natural disasters. It is ranked eighth as the most affected
globally by extreme weather events, and at least seventy-four (74)\% of
Filipinos are vulnerable to multiple hazards. Among these hazards, the most
catastrophic are tropical cyclones, flooding, storm surges, droughts,
earthquakes, and volcanic eruptions. Climate change has caused these events to
increase in intensity, as over the past five years the country has experienced
destructive typhoons. These disasters have caused economic losses and
casualties \parencite{worldbank2023}.

Key challenges in mitigating these risks are the lack of coordination and
insufficient capacity in implementing Disaster Risk Management (DRM) programs,
as most lack integration and cohesiveness to respond to the needs of those
exposed to disasters. Major causes are the lack of execution capacity of
relevant authorities to efficiently mitigate and address adverse events; the
limited coordination between relevant authorities to tackle the increasing
magnitude and frequency of disasters; and the overlapping functions,
competencies, and duplications of roles causing confusion, delays, and
resource misallocations. A core concern is that the timeline of post-disaster
response takes a long time and delays rehabilitation and recovery
\parencite{worldbank2023}.

While social media is utilized for disaster response, the way it is utilized,
and in what context, might point to a more nuanced landscape regarding its
limitations and the nature of its usage by residents in a certain area. A
literature review by \parencite{inferio2025} found that social media disaster
response is not uniformly developed for all aspects, with man-made disasters
not as developed compared to natural disasters in utilization, while COVID-19
and vegetation fires have not been explored as a means for disaster response.

Its utilization was heavy during response and recovery and minimal during preparation. 
Misinformation, noise, and infrastructure limitations are a concern. It is highly effective
in coordinating recovery and providing updates, provided they are filtered
between the informative and the controversial \parencite{inferio2025}---possibly
implying that sophisticated ways to process information and deliver actionable
insights to stakeholders could help communities become more resilient and
responsive in conducting disaster response.

To bridge this gap, this study proposes \emph{SAKLOLO: A Hybrid
	Semantic-Probabilistic Framework for Analysing Social Media Emergency Feeds in
	Bulacan and CAMANAVA}. This framework introduces a novel hybrid architecture
(\sysA) that combines Count Vectorization, Latent Dirichlet Allocation (LDA),
and pre-trained \xlmr-based contextual representations into a Joint Matrix.
This approach preserves mathematical interpretability while achieving high
classification accuracy on chaotic, mixed-syntax Taglish text. Furthermore,
the study conducts a rigorous comparative benchmark of this lightweight system
against a state-of-the-art open-weight dense transformer (\sysB: Qwen3-1.7B\\
-Instruct-Q4\_K\_M) operating under few-shot prompting conditions, 
which is also expected to yield a joint matrix.

The study is focuses on modern-day climate risks and high social media
penetration in the Philippines, as understanding the digital sentiments and
operational needs of vulnerable communities is crucial. This automated
processing allows for a tailored response to specific community needs rather
than a generalized strategy, enabling informed resource allocation of goods,
effort, and time. By benchmarking deterministic and probabilistic machine
learning models as a relative reference against generative AI LLMs, the study ensures that
disaster management offices can focus on their core competencies effective and well aligned 
information processing and dissemination. Rather than manually parsing through an overwhelming
influx of digital noise.

\section*{Statement of the Problem}

This study evaluates and compares the performance of a hybrid joint matrix
pipeline as the baseline (\sysA) against a state-of-the-art large language
model joint matrix output (Qwen3-1.7B\\-Instruct-Q4\_K\_M) in processing unstructured disaster-related
text to generate actionable sentiment classifications and topic modeling for the Bulacan and CAMANAVA regions. Specifically, it seeks to
answer the following questions:

\begin{enumerate}[label=\arabic*.]
	\item Under Bivariate Analysis, to what extent do the topics from the joint matrix of
	both systems show statistically significant association with sentiment
	classifications as measured by:
	\begin{enumerate}[label=\alph*.]
		\item Chi-Square;
		\item Fisher's Exact Test (Contingent on Chi-Square's assumptions not holding).
	\end{enumerate}
	
	\item Is there a statistically significant operational difference between
	the outputs of \sysA{} and \sysB{} as measured
	through:
	\begin{enumerate}[label=\alph*.]
		\item McNemar's Test (with continuity correction).
	\end{enumerate}
\end{enumerate}

\section*{Significance of the Study}

\emph{SAKLOLO: A Hybrid Semantic-Probabilistic Framework for Analysing Social
	Media Emergency Feeds in Bulacan and CAMANAVA} is a computational information
system for disaster response decision-making. It evaluates how unstructured
information in disaster situations can be optimally processed using a Hybrid
LDA + \xlmr{} Joint Matrix compared against a \qwen{} joint matrix
pipeline. By extracting meaningful information---including semantic topics and
contextual sentiment---from unstructured emergency feeds, the study establishes
a mathematically rigorous framework to understand a disaster and respond to it
effectively. This study will benefit several individuals and institutions:

\begin{description}[style=sameline,leftmargin=0pt]
	\item[Bulacan and CAMANAVA Disaster Response Offices.] The study is
	intended to aid disaster response offices in organizing and structuring the
	unorganized disaster-related information they receive through social media
	posts. By determining whether a lightweight hybrid architecture or a massive
	open-weight transformer is more operationally viable for joint latent
	topic modeling and its respective sentiment, DRRMO employees are provided
	with an optimized method to see what disaster-related posts exist and what
	the sentiments regarding these posts are as received by the concerned
	communities. This enables personnel to better understand localized field
	problems and provides them with empirical insight when planning response
	strategies and disseminating information.
	
	\item[Local Government Units (LGUs).] The study assists LGUs in understand-
	ing emergency situations occurring in their specific jurisdictions by
	automatically structuring localized reports. Identifying the most urgent
	issues through sentiment and latent topic analysis helps LGUs gain
	insights into the current state and emotions their respective citizens have
	during natural disasters. This will aid in assessing the disaster situation
	locally and assist LGUs in making decisions that are more aligned, and in
	conducting messaging that takes into consideration how their citizens talk
	about sets of words that have topics, and what their sentiment is
	around those co-occurring sets of words in social media posts.
	
	\item[Emergency Responders.] The study may be beneficial to emergency
	responders by providing disaster-related information in a more organized
	format. As opposed to reviewing social media posts as isolated and
	unorganized streams of information, responders can evaluate co-occurring
	words, their sentiments, and how their respective joint matrices relate to
	each other. This organized intelligence enables emergency responders to make
	highly informed, emotionally intelligent, and community-aligned messaging
	that aids disaster response.
	
	\item[Researchers and Future Developers.] The study acts as a mathematical
	and computational foundation for those interested in natural language
	processing within high-stress, low-resource languages (Taglish).
	Specifically, the novel joint matrix approach combining LDA topic
	probabilities with \xlmr{} contextual embeddings provides reproducible
	evidence on hybrid architecture performance versus dense large language
	models. The integration of Bayesian inference, information theory, and
	comparative deep learning establishes a robust baseline for future researchers
	pursuing text optimization in Philippine disaster management.
\end{description}

\section*{Scope and Delimitation}

This study focuses on the comparative quantitative evaluation of two
computational systems designed to perform topic modeling and sentiment analysis on
unstructured disaster-related reports. \sysA{} is a hybrid pipeline utilizing
Count Vectorization, Latent Dirichlet Allocation (LDA), and \xlmr{}
contextual representations combined into a Joint Matrix 
. \sysB{} utilizes \qwen, an
open-weight large language model, operating under few-shot prompting
conditions, the two systems will be compared in a same measurement comparative analysis.

The geographic scope of the study covers the province of Bulacan and the
CAMANAVA region (Caloocan, Malabon, Navotas, and Valenzuela). The temporal
scope of the corpus spans a five-year period from 2020 to 2025. Data collection
is strictly limited to publicly accessible Facebook posts and comments written
in Tagalog, English, and Taglish that meet the established inclusion criteria
regarding natural disasters (e.g., typhoons, flooding).

The evaluation framework is strictly delimited to specific mathematical and
statistical metrics. Joint Matrix analysis is to be
conducted to analytically test the categorical associations of words to their
sentiment using Chi-square, and a contingent Fisher's Exact Test. The operational
utility of the systems is evaluated through a purposive stakeholder assessment using McNemar's Test.

The study does not involve real-time integration into actual DRRM command
centers, \textbf{nor does it include the physical dispatch of rescue teams}. It performs
social media post backend aggregation with respect to the Data Privacy Act of
2012; all sampled posts and comments are anonymized and protect the identity
of the Facebook users. \sysA{} features a pre-trained \xlmr{}
(\texttt{cardiffnlp/twitter-XLM-RoBERTa-base-sentiment}) that features
approximately one million Tagalog social media posts as part of its training
dataset alongside vast numbers of other languages.

Regarding system implementation, the backend computational pipeline is 
encapsulated within a \textbf{localized analytical web dashboard} developed 
using a lightweight Python-based web framework (e.g., Streamlit). This 
dashboard is strictly delimited to an offline, local-host analytical tool 
designed for the researchers and the expert stakeholder panel. This study does not produce a publicly hosted, 
enterprise-grade web application, nor does it include real-time cloud 
deployment, multi-user authentication, or live integration into actual 
LGU command centers.